\chapter{Introducción}

\section{Objetivo de la práctica}

El objetivo principal de esta práctica es configurar y utilizar \textbf{correo electrónico seguro} mediante el estándar \textbf{S/MIME} (\textit{Secure/Multipurpose Internet Mail Extensions}), incluyendo:

\begin{itemize}[label=\textbullet]
    \item Creación de un certificado digital personal X.509 firmado por una CA raíz
    \item Instalación de certificados en el cliente de correo Mozilla Thunderbird
    \item Intercambio de mensajes firmados y cifrados con otro compañero
    \item Decodificación de mensajes S/MIME cifrados y firmados desde línea de comandos con \texttt{openssl smime}
\end{itemize}

\section{Conceptos fundamentales}

\subsection{Certificados digitales X.509}

Un \textbf{certificado digital} vincula la identidad de un sujeto (nombre, correo electrónico, organización) con su \textbf{clave pública}. Esta vinculación está firmada digitalmente por una \textbf{Autoridad de Certificación (CA)}, lo que permite a cualquier parte que confíe en esa CA verificar la autenticidad del certificado.

La estructura básica de un certificado X.509 incluye:

\begin{itemize}[label=\textbullet]
    \item \textbf{Sujeto (Subject):} Identidad del titular (CN, O, OU, C, Email)
    \item \textbf{Emisor (Issuer):} CA que firma el certificado
    \item \textbf{Clave pública:} La clave pública del titular
    \item \textbf{Periodo de validez:} Fechas de inicio y expiración
    \item \textbf{Firma digital:} Firma de la CA sobre todos los campos anteriores
\end{itemize}

\subsection{S/MIME: Correo electrónico seguro}

S/MIME utiliza criptografía de clave pública y certificados X.509 para proporcionar tres servicios de seguridad al correo electrónico:

\begin{table}[H]
    \centering
    \caption{Serviciones de seguridad de S/MIME}
    \begin{tabularx}{\textwidth}{l|l|X}
        \toprule
        \textbf{Servicio} & \textbf{Mecanismo} & \textbf{Descripción} \\
        \midrule
        Confidencialidad & Cifrado & Solo el destinatario puede leer el mensaje (cifrado con su clave pública) \\
        Autenticación & Firma digital & Se verifica la identidad del emisor \\
        Integridad & Firma digital & Se detecta si el mensaje fue alterado \\
        \bottomrule
    \end{tabularx}
\end{table}

\subsection{Estructura de un mensaje cifrado y firmado}

Cuando un mensaje S/MIME es tanto firmado como cifrado, se crea una estructura de \textbf{doble envolvente}:

\begin{enumerate}
    \item El emisor \textbf{firma} el mensaje con su clave privada $\rightarrow$ genera estructura \texttt{signed-data} (SMIME.P7S)
    \item El resultado se \textbf{cifra} con la clave pública del destinatario $\rightarrow$ genera estructura \texttt{enveloped-data} (SMIME.P7M)
\end{enumerate}

Para decodificarlo, el receptor realiza el proceso inverso:

\begin{enumerate}
    \item \textbf{Descifra} con su clave privada (SMIME.P7M $\rightarrow$ SMIME.P7S)
    \item \textbf{Verifica la firma} con el certificado del emisor (SMIME.P7S $\rightarrow$ texto plano)
\end{enumerate}

\subsection{Formatos de ficheros}

\begin{table}[H]
    \centering
    \caption{Formatos de ficheros utilizados en esta práctica}
    \begin{tabularx}{\textwidth}{l|l|X}
        \toprule
        \textbf{Extensión} & \textbf{Formato} & \textbf{Contenido} \\
        \midrule
        .crt & PEM (Base64) & Certificado X.509 (solo parte pública) \\
        .key & PEM (Base64) & Clave privada (puede estar cifrada con contraseña) \\
        .csr & PEM (Base64) & Solicitud de firma de certificado \\
        .p12 & PKCS\#12 (binario) & Certificado + clave privada empaquetados \\
        .eml & Texto & Mensaje de correo completo (cabeceras + cuerpo MIME) \\
        .p7m & PKCS\#7 (binario) & Mensaje cifrado S/MIME (\texttt{enveloped-data}) \\
        .p7s & PKCS\#7 (binario) & Mensaje firmado S/MIME (\texttt{signed-data}) \\
        \bottomrule
    \end{tabularx}
\end{table}

\section{Herramientas utilizadas}

\begin{itemize}[label=\textbullet]
    \item \textbf{OpenSSL:} Herramienta de criptografía en línea de comandos
    \item Comandos principales: \texttt{genpkey}, \texttt{req}, \texttt{x509}, \texttt{pkcs12}, \texttt{smime}
    \item \textbf{Mozilla Thunderbird:} Cliente de correo electrónico con soporte S/MIME
\end{itemize}

